\section{Exhibit B-1: prime closure diagnostics}
\label{sec:exhibit-prime}

This exhibit applies the \SBT closure lens to a canonical ``discrete $\rightarrow$ spectral'' packaging problem:
compressing the multiplicative micro-structure of the integers (primes) into a continuous analytic object.
We do \emph{not} attempt to prove theorems about $\zeta$.
Instead we build falsification-first diagnostics that separate:
(i) a convergence control regime where the micro-descriptions are valid, from
(ii) the critical strip where naive staging is unstable and protocol mismatch blows up.

\subsection{Two staged micro-descriptions}
Fix $s\in\mathbb{C}$.
The classical zeta function has two micro-descriptions in a common domain (and, in the infinite limit, they agree there):
\[
\zeta(s)\;=\;\sum_{n\ge 1} n^{-s}
\qquad\text{and}\qquad
\zeta(s)\;=\;\prod_{p}\frac{1}{1-p^{-s}},
\quad (\Re(s)>1).
\]
These classical representations and their domains are standard; see e.g.\ \cite{Titchmarsh1986}.
At finite stage we replace these by staged truncations.
In the recorded diagnostics we use a mild exponential smoothing weight $w(u)=e^{-u}$ and define:
\begin{align*}
S_N(s) &:= \sum_{n=1}^{N} w(n/N)\,n^{-s},
\\
P_N(s) &:= \exp\!\left(\sum_{p\le N} w(p/N)\,\bigl(-\log(1-p^{-s})\bigr)\right),
\end{align*}
where the Euler-side expression is computed via a weighted log-sum for numerical stability.
Note that the smoothing is applied at the level of the log-sum:
equivalently,
\[
P_N(s) \;=\; \prod_{p\le N} (1-p^{-s})^{-w(p/N)}.
\]
We use this only to reduce sharp cutoff artifacts; when $w\equiv 1$ this reduces to the usual sharp partial Euler product.
The key \SBT point is conceptual:
$S_N$ is an \emph{additive} staging protocol and $P_N$ is a \emph{multiplicative} staging protocol.
A stable closure would require these protocols to become compatible under refinement \Pfour\ and under the chosen packaging \Pfive, as audited by an accounting ledger \Psix.

\subsection{Prototype packaging maps: raw, completed, symmetrized}
To compare staged objects across regions, we apply simple prototype packaging maps.
Let
\[
C(s)\;:=\;\tfrac12\,s(s-1)\,\pi^{-s/2}\Gamma(s/2).
\]
Given a staged function $F(s)$, define:
\begin{align*}
\mathrm{Raw}(F)(s) &:= F(s),
\\
\mathrm{Comp}(F)(s) &:= C(s)\,F(s),
\\
\mathrm{Sym}(F)(s) &:= \tfrac12\bigl(\mathrm{Comp}(F)(s)+\mathrm{Comp}(F)(1-s)\bigr).
\end{align*}
Here $\mathrm{Comp}$ is a one-sided ``completion prototype'', and $\mathrm{Sym}$ enforces the self-dual symmetry $s\leftrightarrow 1-s$ by construction.
We emphasize that this is a \emph{diagnostic} packaging, not an analytic continuation scheme.

In the recorded run we choose:
\begin{itemize}
  \item \textbf{Convergence region:} $\mathrm{Comp}$ (one-sided completion), because it does not force evaluation at $1-s$ far outside the domain where these truncations represent $\zeta$ (for $\Re(s)>1$ one has $\Re(1-s)<0$).
  \item \textbf{Critical strip:} $\mathrm{Sym}$ (symmetrized completion), as a prototype of ``self-dual packaging''.
\end{itemize}

\subsection{Diagnostics: route mismatch and error ledger}
Let $\mathrm{Pack}$ denote the chosen packaging mode in a given regime.
On a fixed finite test set $K$ (a small grid of points in the region of interest), we evaluate the packaged staged functions.
In the recorded run, $K$ is a grid with $\sigma\in$ \PrimeSigmaConvList\ and $t\in$ \PrimeTList\ in the convergence regime, and $\sigma\in$ \PrimeSigmaStripList\ and $t\in$ \PrimeTList\ in the strip; refinements use $N\in$ \PrimeNList.
\[
A_N(s):=\mathrm{Pack}(S_N)(s),
\qquad
B_N(s):=\mathrm{Pack}(P_N)(s).
\]
We quantify protocol mismatch by a normalized route mismatch:
\[
\mathrm{RM}_2(N)
\;:=\;
\frac{\bigl(\sum_{s\in K}|A_N(s)-B_N(s)|^2\bigr)^{1/2}}
     {\bigl(\sum_{s\in K}|A_N(s)|^2\bigr)^{1/2}+\varepsilon_0},
\]
with a small stabilizer $\varepsilon_0>0$.

All reported values in this exhibit are computed at working precision \PrimeDPS\ digits (via \texttt{mpmath}).

Separately, we compute an \Psix\ error ledger by comparing each packaged truncation to a reference value computed by \texttt{mpmath.zeta} under the same packaging mode:
\[
\mathrm{Err}_S(N) \sim \frac{\|A_N-\mathrm{Pack}(\zeta)\|}{\|\mathrm{Pack}(\zeta)\|+\varepsilon_0},
\qquad
\mathrm{Err}_P(N) \sim \frac{\|B_N-\mathrm{Pack}(\zeta)\|}{\|\mathrm{Pack}(\zeta)\|+\varepsilon_0}.
\]
In the convergence control region this is a genuine accuracy test.
In the critical strip it is a stress test: it measures how badly naive staging fails to track the analytically continued object.

\subsection{Control regime: \texorpdfstring{$\Re(s)>1$}{Re(s)>1}}
In the convergence control regime (where both series and product are valid in the infinite limit), the diagnostics behave sensibly:
both the mismatch $\mathrm{RM}_2(N)$ and the errors decrease with $N$ on the recorded test set.
Table~\ref{tab:prime-rm-conv} shows $\mathrm{RM}_2(N)$ together with the two error ledgers.
In particular, the route mismatch decreases to $\mathrm{RM}_2(800)=$ \PrimeRMConvAtEightHundred,
consistent with staged compatibility improving under refinement when the micro-descriptions are valid.

% Auto-generated by scripts/export_results_tex.py from notes/*_last_run.json
\begin{table}[h]
\centering
\caption{Prime RM in convergence region}
\label{tab:prime-rm-conv}
\begin{tabular}{l|l|l|l}
\hline
 N & RM2 & errS2 & errP2 \\
\hline
 50 & \ensuremath{0.0759745} & \ensuremath{0.209083} & \ensuremath{0.146987} \\
\hline
 100 & \ensuremath{0.0697803} & \ensuremath{0.171755} & \ensuremath{0.111871} \\
\hline
 200 & \ensuremath{0.0667314} & \ensuremath{0.142319} & \ensuremath{0.084813} \\
\hline
 400 & \ensuremath{0.0613035} & \ensuremath{0.118486} & \ensuremath{0.0654121} \\
\hline
 800 & \ensuremath{0.0523991} & \ensuremath{0.0984851} & \ensuremath{0.0505677} \\
\hline
\end{tabular}
\end{table}


\subsection{Critical strip: protocol mismatch under invalid staging}
In the critical strip ($0<\Re(s)<1$), the naive truncations $S_N$ and especially $P_N$ are \emph{not} valid micro-descriptions of the analytically continued $\zeta$.
Accordingly, the diagnostics indicate a non-feasible closure:
the route mismatch and errors grow rapidly with $N$ on the recorded test set.
Table~\ref{tab:prime-rm-strip} shows that the mismatch reaches $\mathrm{RM}_2(800)=$ \PrimeRMStripAtEightHundred,
and the Euler-side error ledger becomes enormous.

% Auto-generated by scripts/export_results_tex.py from notes/*_last_run.json
\begin{table}[h]
\centering
\caption{Prime RM in the critical strip}
\label{tab:prime-rm-strip}
\begin{tabular}{l|l|l|l}
\hline
 N & RM2 & errS2 & errP2 \\
\hline
 50 & \ensuremath{340.421} & \ensuremath{6.59171} & \ensuremath{2123.75} \\
\hline
 100 & \ensuremath{5996.54} & \ensuremath{10.4329} & \ensuremath{6.02608\times 10^{4}} \\
\hline
 200 & \ensuremath{6.19843\times 10^{5}} & \ensuremath{16.5012} & \ensuremath{0.984885\times 10^{7}} \\
\hline
 400 & \ensuremath{3.14454\times 10^{8}} & \ensuremath{27.1202} & \ensuremath{8.29135\times 10^{9}} \\
\hline
 800 & \ensuremath{5.00808\times 10^{12}} & \ensuremath{43.7077} & \ensuremath{2.16627\times 10^{14}} \\
\hline
\end{tabular}
\end{table}


\begin{figure}[htbp]
  \centering
  \includegraphics[width=0.85\linewidth]{prime_closure_rm_last_run.png}
  \caption{Prime closure route mismatch diagnostic. The same staged additive and multiplicative protocols behave compatibly in a convergence control regime (Re(s)>1), but mismatch grows by many orders of magnitude in the critical strip under naive staging and prototype self-dual packaging. This is interpreted as a feasibility failure of the staging/packaging pair, not as a theorem about $\zeta$.}
  \label{fig:prime-rm}
\end{figure}

\subsection{Interpretation: a diagnostic, not a theorem about \texorpdfstring{$\zeta$}{zeta}}
This exhibit is deliberately modest.
It does not claim that $\mathrm{Sym}$ is ``the right'' packaging, nor that the observed blow-up has direct implications for the zero set of $\zeta$.
Instead it instantiates the \SBT logic:

\begin{itemize}
  \item \textbf{Multiple routes exist} (additive vs multiplicative staging), hence \Pthree\ applies.
  \item \textbf{Feasibility must be audited} by a ledger (\Psix): mismatch and error should shrink in stable regimes.
  \item \textbf{Staging matters} (\Pfour): outside the domain of convergence, naive truncations are not compatible with the packaged object one hopes to define.
\end{itemize}

The positive outcome is methodological: the same route-mismatch machinery used for discrete derivatives (Exhibits~A-1 and A-2) can be used to stress-test prime-to-spectral closure claims.
A stable analytic continuation scheme---if expressed as a refinement-fixed packaging with an appropriate positivity/passivity ledger---should manifest as decreasing mismatch under refinement on suitable diagnostics.
Here we only establish the baseline: without additional feasibility structure, the mismatch blows up exactly where classical analytic number theory says naive micro-descriptions fail.
